% !TeX program = pdflatex
% Standard single-column preprint; no custom class, style, or external fonts.
% Formatting follows https://info.arxiv.org/help/submit_tex.html.
% TODO: Supply author names, affiliations, and contact details before submission.
% The blank author field is not an anonymous-review declaration.
% AI assistance was used for translation, technical revision, and reference-code preparation.
% Author review of the revised manuscript and reference implementation remains pending.
% The new reference run is documented in paper/reproducibility/README.md.
% Its recorded results replace the earlier draft's unavailable reference run.
\documentclass[11pt,a4paper]{article}
\usepackage[T1]{fontenc}
\usepackage[utf8]{inputenc}
\usepackage[english]{babel}
\usepackage{lmodern}
\usepackage[margin=1in]{geometry}
\usepackage{amsmath,amssymb,amsthm,bm,mathtools}
\usepackage{microtype}
\usepackage{booktabs,tabularx,array}
\usepackage{algorithm,algpseudocode}
\usepackage[numbers,sort&compress]{natbib}
\usepackage{enumitem}
\usepackage[hidelinks,bookmarksnumbered=true]{hyperref}
\usepackage{xurl}
\hypersetup{
  pdftitle={A Graph-Tensor Operator Formulation of Classical Molecular Dynamics: Equivalent Reformulations and Hardware-Independent Execution Semantics},
  pdfauthor={},
  pdfsubject={Preprint: mathematical formulation and small-system equivalence checks},
  pdfkeywords={molecular dynamics, graph operators, adjoint assembly, many-body potentials}
}
\setlength{\tabcolsep}{5pt}
\renewcommand{\arraystretch}{1.12}
\setlist{nosep,leftmargin=2em,topsep=.25em}
\setlength{\emergencystretch}{2em}
\allowdisplaybreaks[2]
\raggedbottom
\setcounter{secnumdepth}{3}
\pagestyle{plain}
\algrenewcommand\algorithmicrequire{\textbf{Input:}}
\algrenewcommand\algorithmicensure{\textbf{Output:}}
\newcommand{\oldform}{\par\noindent\textbf{Conventional form.}\quad}
\newcommand{\newform}{\par\noindent\textbf{Operator form.}\quad}
\newcommand{\parallelmeaning}{\par\noindent\textbf{Equivalence.}\quad}
\newcommand{\RR}{\mathbb R}
\newcommand{\NN}{\mathbb N}
\newcommand{\ZZ}{\mathbb Z}
\newcommand{\ones}{\mathbf 1}
\newcommand{\dd}{\mathrm d}
\newcommand{\tr}{\operatorname{tr}}
\newcommand{\diag}{\operatorname{diag}}
\newcommand{\Map}{\operatorname{Map}}
\newcommand{\Sum}{\operatorname{Sum}}
\newcommand{\Gather}{\operatorname{Gather}}
\newcommand{\Assemble}{\operatorname{AssembleSum}}
\newcommand{\SegSum}{\operatorname{SegmentSum}}
\newcommand{\atan}{\operatorname{atan2}}
\newcommand{\round}{\operatorname{round}}
\newcommand{\argmin}{\operatorname*{arg\,min}}
\newcommand{\ip}[2]{\left\langle #1,#2\right\rangle}
\newcommand{\norm}[1]{\left\lVert #1\right\rVert}
\newcommand{\vect}[1]{\bm{#1}}
\newcommand{\mass}{\mathsf M}
\newcommand{\cell}{\mathsf H}
\newcommand{\src}{S_{\mathrm s}}
\newcommand{\dst}{S_{\mathrm t}}
\newcommand{\GE}{\mathcal G_I}
\newtheorem{proposition}{Proposition}[section]
\theoremstyle{definition}
\newtheorem{definition}[proposition]{Definition}
\newtheorem{remark}[proposition]{Remark}
\newcolumntype{Y}{>{\raggedright\arraybackslash}X}
\newcommand{\paperauthors}{} % Author Name\\Affiliation\\\texttt{email@example.org}
\title{A Graph--Tensor Operator Formulation of\\
Classical Molecular Dynamics:\\[0.4em]
\large Equivalent Reformulations and Hardware-Independent Execution Semantics}
\author{\paperauthors}
\date{September 17, 2026}

\begin{document}
\maketitle
\thispagestyle{plain}
\begin{abstract}
Interaction sums in classical molecular dynamics are inherently data parallel, yet interaction
indices, geometric transformations, local differentiation, and force accumulation are often
intertwined with implementation details. We present a hardware-independent graph--tensor
operator formulation that recasts conventional indexed expressions as compositions of local
extraction, geometric maps, energy transformations, and adjoint assembly while preserving the
potential energy, cutoff conventions, and equations of motion. A directed incidence operator
yields a unified expression for pair forces. The formulation extends to bonds, angles, and
dihedrals on ordered hyperedges, and to embedded-atom models with node--edge aggregation. We
specify conditions for periodic images and neighbor-list validity, express velocity Verlet
integration as a composition of state maps, and describe representation invariance and
continuous conservation properties. Double-precision tests on small systems examine energy
gradients, scalar-loop and operator agreement, half- and full-list assembly, the
extraction--assembly adjoint identity, and single-step integration equivalence.
The results support mathematical consistency for the tested
examples; they do not validate GPU performance or a complete simulation engine. The formulation
separates physical semantics from data layout, operator fusion, and hardware scheduling,
providing testable mathematical interfaces for subsequent implementations.
\end{abstract}
\noindent\textbf{Keywords:} classical molecular dynamics; sparse interaction graphs; adjoint assembly; many-body potentials; hardware-independent operators

\section{Introduction}
\label{sec:intro}
Molecular dynamics (MD) evolves a particle system using a specified potential energy and its
gradient. For a given physical model, atom-wise loops, interaction-wise evaluations, and
reductions on a graph can represent the same energy and forces. Their primary distinction lies
in the organization of data dependencies. Existing parallel simulation systems already make
extensive use of spatial decomposition, neighbor lists, and accelerator
execution~\cite{thompson2022,anderson2020}. Rewriting indexed sums as tensor operations should
therefore not be interpreted as introducing parallelism into MD for the first time.

We address a more specific representational question: can interaction indexing, local geometry,
chain-rule differentiation, and global accumulation each be assigned an explicit mathematical
interface while retaining testable equivalence across implementations? We express interaction
extraction as a linear operator, retain local geometry as a nonlinear map, and define force
accumulation through the adjoint of extraction. Mathematical intermediate tensors need not
correspond to separate buffers: implementations remain free to choose storage layouts, fusion
boundaries, and reduction strategies.

The paper develops this formulation in three parts. It establishes a unified relation between
extraction, assembly, and local energy gradients; derives conventional and operator expressions
for pair potentials, bonded interactions, embedded-atom models, and time integration; and checks
signs, counting factors, and chain-rule consistency through reproducible small-system tests. The
subject is a computational formulation and its contracts, rather than a new potential-energy
model. We use established graph, tensor, and differentiation tools to develop a consistent
representation of the classical MD models considered here, without presupposing a performance
advantage over existing software.

\subsection{Related work and scope}
LAMMPS provides a parallel simulation framework for particle models spanning multiple
scales~\cite{thompson2022}. Kokkos supports portable implementations on multicore processors and
accelerators through data-layout and parallel-execution abstractions~\cite{edwards2014}; such
abstractions are compatible with the mathematical formulation developed here and could serve as
implementation backends. HOOMD-blue organizes particle simulations and their data structures
around GPU acceleration~\cite{anderson2020}. JAX-MD constructs simulations from differentiable
functions, array-based states, and spatial partitioning operations, and supports neural-network
interactions~\cite{schoenholz2020}. These systems already encompass many central ideas in GPU
execution and array-based MD. More specifically, JAX-MD's structure-mapping functions lift
pair and bonded functions to collections of interactions and support neighbor-list
evaluation~\cite{jax-smap}. The extraction, local evaluation, and reduction pattern is therefore
an established computational construction.

The physical ingredients also have established precedents: the EAM formulation of Daw and
Baskes~\cite{daw1984}, torsion-angle derivatives studied by Blondel and
Karplus~\cite{blondel1996}, reversible operator splitting by Tuckerman, Berne, and
Martyna~\cite{tuckerman1992}, and periodic many-body virial formulations by Thompson,
Plimpton, and Mattson~\cite{thompson2009}. We do not claim these ingredients or the
Gather--assembly adjoint rule as new. The contribution is their organization into a common set
of explicit mathematical contracts, with aligned sign, counting, geometry, and differentiation
conventions and executable checks for the examples considered here.

In contrast to these complete software systems, this paper is limited to a methodological
formulation: explicit transformations from indexed equations to graph--tensor operators, reverse
assembly rules, correctness conditions, and a restricted set of reference checks. Applicability
and performance must be established by subsequent backend implementations and independent
benchmarks.
\section{Problem definition and notation}
\label{sec:notation}
\subsection{State and index spaces}
Let the position, velocity, and force of particle $i$ be column vectors $\vect x_i,\vect
v_i,\vect f_i\in\RR^3$. Their transposes form row $i$ of the corresponding arrays:
\begin{equation}
X_{i,:}=\vect x_i^{\mathsf T},\qquad
V_{i,:}=\vect v_i^{\mathsf T},\qquad
F_{i,:}=\vect f_i^{\mathsf T}.
\end{equation}
The gradient $\nabla_X U$ has the same shape as $X$, with the differential defined using the
Frobenius inner product:
\begin{equation}
\dd U=\ip{\nabla_XU}{\dd X}_F
=\sum_{i=1}^N\sum_{\alpha=1}^3
\frac{\partial U}{\partial X_{i\alpha}}\,\dd X_{i\alpha}.
\end{equation}

\begin{table}[htbp]
\centering
\caption{Principal symbols and array shapes. Three-dimensional particle vectors are stored as rows. Indices, parameters, and continuous state variables have distinct mathematical roles.}
\label{tab:notation}
\small
\begin{tabularx}{\textwidth}{@{}p{22mm}p{45mm}Y@{}}
\toprule
Symbol & Shape or definition & Meaning\\
\midrule
$X,V,F$ & $\RR^{N\times3}$ & Positions, velocities, and forces; continuous variables\\
$m,\mass$ & $m\in\RR_{>0}^N$, $\mass=\diag(m)$ & Atomic masses; $\mass^{-1}F$ divides each row by its mass\\
$z,\Theta$ & $z\in\NN^N$; parameter set & Element/type labels and fixed potential parameters\\
$U,K,\mathcal H$ & Scalars & Potential, kinetic, and total energy; $\mathcal H=K+U$\\
$P$ & Nonnegative integer & Number of unordered pair interactions, distinct from energy\\
$I^{(k)}$ & $\{1,\ldots,N\}^{P_k\times k}$ & Ordered atom indices for each $k$-body term\\
$B,C$ & $\RR^{P\times N}$ & Directed incidence and unsigned endpoint operators\\
$D,r$ & $\RR^{P\times3}$, $\RR_{>0}^{P}$ & Edge displacements and distances\\
$\cell,\Omega$ & $\RR^{3\times3}$, $|\det\cell|$ & Cell matrix (lattice vectors as rows) and volume\\
\bottomrule
\end{tabularx}
\end{table}

Table~\ref{tab:notation} summarizes the notation. Element labels, masses, and parameters may be
stored in arrays without becoming a common class of ``learnable features.'' In particular,
indices define discrete structure, whereas positions are continuous variables. Fixed parameters
and discrete indices are not differentiated as position components when computing position
gradients.

\subsection{Assumptions}
The main derivations assume a fixed number of particles, fixed positive masses, a conservative
potential energy, and unconstrained motion. Unless periodic boundaries are specified, the
equations apply in nonperiodic Euclidean space. Pair edges $(i_e,j_e)$ satisfy $i_e\ne j_e$,
with each unordered atom pair appearing once; edge direction is a notational choice. Symmetric
pair interactions satisfy $\phi_{ij}=\phi_{ji}$.

Periodic nonbonded examples use an invertible, fixed cell and one closest image for each
distinct unordered atom pair, with no self-image terms. This defines a minimum-image model;
it is not, in general, a sum over all lattice images. Section~\ref{sec:pbc} states a sufficient
cutoff condition under which the two short-range descriptions agree.

For periodic systems, the matrix expressions are differentiated \textbf{locally with the current
neighbor relation and image branch held fixed}, away from image ties and other nonsmooth
points of the specified potential. No cutoff is needed for this local chain rule.
Image selection and neighbor validity must still be updated during a simulation, as discussed
in Section~\ref{sec:pbc}.


\paragraph{Physical semantics and implementation freedom.}
The specified potential-energy model, cutoff conventions, units, exclusions, and equations of
motion remain unchanged.

Index organization, edge/atom grouping, storage layout, summation order, operator fusion, and
hardware scheduling may change. Mathematical equivalence holds in real arithmetic;
floating-point implementations additionally require error and reproducibility specifications.



\subsection{Conventional dynamics and tensor state}
\label{sec:baseline}
\oldform Newton's equations and the conservative-force relation apply to each atom:
\begin{equation}
\frac{\dd\vect x_i}{\dd t}=\vect v_i,\qquad
m_i\frac{\dd\vect v_i}{\dd t}=\vect f_i,\qquad
\vect f_i=-\frac{\partial U}{\partial\vect x_i}.
\label{eq:newton-classic}
\end{equation}
The kinetic and total energies are
\begin{equation}
K=\frac12\sum_{i=1}^N m_i\norm{\vect v_i}^2,
\qquad \mathcal H=K+U.
\end{equation}
These relations do not depend on the particular potential or processor
architecture~\cite{gromacs-md}.

\newform The same relations can be written using arrays and operators:
\begin{equation}
\dot X=V,\qquad \dot V=\mass^{-1}F,\qquad F=-\nabla_XU(X;\Theta).
\label{eq:newton-tensor}
\end{equation}
\begin{equation}
K=\frac12\tr(V^{\mathsf T}\mass V).
\end{equation}
This step unifies the state representation without yet changing the organization of force
evaluation.

\section{Extraction, adjoint assembly, and local energy gradients}
\label{sec:gather}
\subsection{Conventional local interactions}
For a class of $k$-body interactions, let term $a$ involve atoms
$(I_{a,1},\ldots,I_{a,k})$. The conventional energy is
\begin{equation}
U_I(X)=\sum_{a=1}^{P_k}
u_a\bigl(\vect x_{I_{a,1}},\ldots,\vect x_{I_{a,k}};\Theta_a\bigr).
\label{eq:local-classic}
\end{equation}
The gradient for each atom sums the local gradients of all interactions involving that atom,
giving the force
\begin{equation}
\vect f_i=-\sum_{\substack{a,p\\I_{a,p}=i}}
\nabla_p u_a\bigl(\vect x_{I_{a,1}},\ldots,\vect x_{I_{a,k}};\Theta_a\bigr).
\label{eq:local-classic-force}
\end{equation}
Here $\nabla_p$ means the gradient with respect to the $p$th argument before identifying
arguments with global atom coordinates. If an atom occupies several slots, each slot contributes.
Roles and ordering within an interaction must be preserved; for example, the second atom in an
angle is its central atom.

\subsection{Gather and its adjoint}
\begin{definition}[Local extraction]
For fixed indices $I$, define the linear operator
\begin{equation}
\GE:\RR^{N\times3}\to\RR^{P_k\times k\times3},\qquad
[\GE(X)]_{a,p,\alpha}=X_{I_{a,p},\alpha}.
\label{eq:gather}
\end{equation}
\end{definition}
\begin{definition}[Additive assembly]
Define
\begin{equation}
\GE^*:\RR^{P_k\times k\times3}\to\RR^{N\times3},\qquad
[\GE^*(Z)]_{i,\alpha}=\sum_{\substack{a,p\\I_{a,p}=i}} Z_{a,p,\alpha}.
\label{eq:assemble}
\end{equation}
\end{definition}
Contributions at repeated indices must be summed, rather than overwritten or averaged.

\begin{proposition}[Extraction--assembly adjoint identity]
\begin{equation}
\ip{\GE(X)}{Z}=\ip{X}{\GE^*(Z)}.
\label{eq:adjoint-identity}
\end{equation}
\end{proposition}
\begin{proof}
Regrouping the indexed sum gives
\begin{align}
\ip{\GE(X)}{Z}
&=\sum_{a,p,\alpha}X_{I_{a,p},\alpha}Z_{a,p,\alpha}\\
&=\sum_{i,\alpha}X_{i,\alpha}
\sum_{a,p:I_{a,p}=i}Z_{a,p,\alpha}
=\ip{X}{\GE^*(Z)}.\nonumber
\end{align}
Thus $\GE^*$ is the adjoint of $\GE$ under the chosen inner products.
\end{proof}

\subsection{Reformulating energy and forces}
Let $Q=\GE(X)$ and $\Phi_I(Q)=\sum_a u_a(Q_a)$. Then
\begin{equation}
U_I(X)=\Phi_I(\GE(X)),\qquad
F_I=-\GE^*\bigl(\nabla_Q\Phi_I(Q)\bigr).
\label{eq:local-reform}
\end{equation}
The result follows directly from the chain rule:
\begin{equation}
\dd U_I=\ip{\nabla_Q\Phi_I}{\GE(\dd X)}
=\ip{\GE^*\nabla_Q\Phi_I}{\dd X}.
\end{equation}

If local coordinates are first mapped to distances, angles, or other geometric quantities
$Y=\Gamma_I(Q)$, the expression becomes
\begin{equation}
U_I=\Phi_I\circ\Gamma_I\circ\GE(X),\qquad
F_I=-\GE^*\,\mathrm D\Gamma_I(Q)^*\,
\nabla_Y\Phi_I(Y).
\label{eq:geometry-chain}
\end{equation}
Here $\mathrm D\Gamma_I(Q)^*$ denotes the transpose action of the geometric Jacobian. The linear
Gather operator has a fixed adjoint, whereas the derivative of nonlinear geometry depends on the
current coordinates.


\paragraph{Forward and reverse dependencies.}
The forward chain is local extraction $\to$ geometry $\to$ local energy $\to$ summation.

The reverse chain is energy derivatives $\to$ geometric Jacobian transpose action $\to$ additive
assembly $\to$ negation to obtain forces.

These are mathematical dependencies. They do not require materialized intermediate arrays or the
use of automatic-differentiation software.


\section{Equivalent reformulations of representative interactions}
\label{sec:models}
\subsection{Pair potentials and graph incidence operators}
\label{sec:pair}

Let $\vect d_{ij}=\vect x_j-\vect x_i$, $r_{ij}=\norm{\vect d_{ij}}$, and let $\mathcal N(i)$ be
a symmetric neighbor set. Counting each unordered interaction once gives
\begin{equation}
U_2=\sum_{i<j}\phi_{ij}(r_{ij})
=\frac12\sum_i\sum_{j\in\mathcal N(i)}\phi_{ij}(r_{ij}).
\label{eq:pair-traditional-energy}
\end{equation}
The sums include only the specified active interactions. The corresponding atomic force is
\begin{equation}
\vect f_i=\sum_{j\in\mathcal N(i)}
\phi'_{ij}(r_{ij})\frac{\vect x_j-\vect x_i}{r_{ij}}.
\label{eq:pair-traditional-force}
\end{equation}
The sign follows from the displacement convention $i\to j$: the expression has a positive
prefactor, while $\phi'(r)$ itself may be positive or negative.


\subsubsection{Selection operators from edge indices}
For $P$ edges $e=(i_e,j_e)$, define
\begin{equation}
(\src)_{ei}=\begin{cases}1&i=i_e,\\0&\text{otherwise},\end{cases}
\qquad
(\dst)_{ei}=\begin{cases}1&i=j_e,\\0&\text{otherwise}.\end{cases}
\end{equation}
Both operators have shape $P\times N$. Set
\begin{equation}
B=\dst-\src,\qquad C=\src+\dst=|B|.
\label{eq:incidence}
\end{equation}
Each row of $B$ has exactly two nonzero entries, $-1$ and $+1$, while $C$ has a $1$ at both
endpoints. By definition,
\begin{equation}
D=BX=(\dst X)-(\src X),\qquad D_e=(\vect x_{j_e}-\vect x_{i_e})^{\mathsf T}.
\label{eq:edge-difference}
\end{equation}

\subsubsection{From conventional energy to edge-wise potentials}
\oldform
\begin{equation}
U_2=\sum_{\{i,j\}\in\mathcal E}\phi_{ij}(\norm{\vect x_j-\vect x_i}).
\end{equation}
\newform With $\phi(r)_e=\phi_e(r_e;\theta_e)$ and $r_e=\sqrt{\sum_\alpha D_{e\alpha}^2}$,
\begin{equation}
U_2(X)=\ones_P^{\mathsf T}\phi(r),\qquad r=\operatorname{RowNorm}(BX).
\label{eq:pair-reform-energy}
\end{equation}
The parameter $\theta_e$ may be obtained by a lookup using the type pair $z_{i_e},z_{j_e}$; it
need not be physically replicated on every edge.

\subsubsection{Energy differentials and adjoint force assembly}
For $r_e>0$,
\begin{equation}
\dd r_e=\frac{D_e\cdot\dd D_e}{r_e},\qquad
g_e=\frac{\partial U_2}{\partial D_e}
=\frac{\phi'_e(r_e)}{r_e}D_e.
\end{equation}
Consequently,
\begin{align}
\dd U_2
&=\sum_e\phi'_e(r_e)\,\dd r_e
=\sum_e g_e\cdot\dd D_e\\
&=\ip{g}{B\,\dd X}_F
=\ip{B^{\mathsf T}g}{\dd X}_F.\nonumber
\end{align}
Comparing with the gradient definition yields
\begin{equation}
F=-B^{\mathsf T}g
=-B^{\mathsf T}\diag\!\left(\frac{\phi'(r)}{r}\right)D.
\label{eq:pair-reform-force}
\end{equation}
For edge $e=(i,j)$, assembly is exactly
\begin{equation}
\vect f_i\mathrel{+}=g_e^{\mathsf T},\qquad
\vect f_j\mathrel{-}=g_e^{\mathsf T}.
\label{eq:pair-endpoints}
\end{equation}
This agrees with Eq.~\eqref{eq:pair-traditional-force}.

\subsubsection{Nonlinearity and sparse implementation}
In the absence of periodic image offsets, one may further write
\begin{equation}
F=-L_w(X)X,\qquad
L_w(X)=B^{\mathsf T}\diag(w(X))B,\qquad
w_e(X)=\frac{\phi'_e(r_e(X))}{r_e(X)}.
\label{eq:nonlinear-laplacian}
\end{equation}
Since $L_w$ depends on $X$, this is generally not a fixed linear system. The weights $w_e$ may
also be negative, so the operator cannot unconditionally be treated as a positive-semidefinite graph
Laplacian.

An implementation needs only $i_e,j_e$. The action $BX$ can be evaluated by indexed reads and
subtraction, and $B^{\mathsf T}g$ by endpoint accumulation or grouped reduction. A backend may
fuse these operations without explicitly storing $B$, $L_w$, or all of $D$.

\subsubsection{A direct check of translation invariance}
For edges $(1,2),(2,3)$,
\begin{equation}
B=\begin{bmatrix}-1&1&0\\0&-1&1\end{bmatrix},\qquad
F=\begin{bmatrix}g_1\\-g_1+g_2\\-g_2\end{bmatrix}.
\end{equation}
Hence $\sum_i\vect f_i=0$. More generally, $B\ones_N=0$, so
\begin{equation}
\ones_N^{\mathsf T}F=-(B\ones_N)^{\mathsf T}g=0.
\label{eq:net-force}
\end{equation}

\subsubsection{Counting equivalence of half and full lists}
LAMMPS distinguishes a half list, which stores each pair once, from a full list, which stores
both directions~\cite{lammps-neighbor}. We use half lists unless stated otherwise.

\oldform For a complete bidirectional neighbor set,
\begin{equation}
U_2=\frac12\sum_i\sum_{j\in\mathcal N(i)}\phi_{ij}(r_{ij}),\qquad
\vect f_i=\sum_{j\in\mathcal N(i)}g_{i\to j}^{\mathsf T}.
\end{equation}
\newform If $\widetilde B,\widetilde\src,\widetilde g$ are defined on the complete set of directed edges,
\begin{equation}
U_2=\frac12\ones^{\mathsf T}\phi(\widetilde r),\qquad
F=-\frac12\widetilde B^{\mathsf T}\widetilde g
=\widetilde\src^{\mathsf T}\widetilde g.
\label{eq:full-list}
\end{equation}
The final equality requires all reverse edges and a symmetric potential, so that $g_{j\to
i}=-g_{i\to j}$.



For a half list, each unordered pair is evaluated once. The energy has no factor of $1/2$, and
forces are accumulated at both endpoints.

For a full list, the energy carries a factor of $1/2$. If each atom collects only the forces
from its outgoing edges, that force sum \textbf{has no additional factor of $1/2$}.
Differentiating the full-list energy through the complete chain instead gives
$-\tfrac12\widetilde B^{\mathsf T}\widetilde g$.


\subsection{The Lennard--Jones potential and cutoff conventions}
\label{sec:lj}
\oldform The standard 12--6 Lennard--Jones (LJ) potential is~\cite{lammps-lj}
\begin{equation}
\phi_{ij}(r)=4\epsilon_{ij}\left[
\left(\frac{\sigma_{ij}}r\right)^{12}
-\left(\frac{\sigma_{ij}}r\right)^6\right].
\label{eq:lj-energy}
\end{equation}
Direct differentiation gives
\begin{equation}
\phi'_{ij}(r)=\frac{24\epsilon_{ij}}r
\left[\left(\frac{\sigma_{ij}}r\right)^6
-2\left(\frac{\sigma_{ij}}r\right)^{12}\right].
\end{equation}
The atomic force is therefore
\begin{equation}
\vect f_i=\sum_j\frac{24\epsilon_{ij}}{r_{ij}^2}
\left[\left(\frac{\sigma_{ij}}{r_{ij}}\right)^6
-2\left(\frac{\sigma_{ij}}{r_{ij}}\right)^{12}\right]
(\vect x_j-\vect x_i).
\label{eq:lj-force-traditional}
\end{equation}

\newform Define the edge-wise quantities
\begin{equation}
a_e=\left(\frac{\sigma_e}{r_e}\right)^6,\qquad
u_e=4\epsilon_e(a_e^2-a_e),\qquad
w_e=\frac{24\epsilon_e}{r_e^2}(a_e-2a_e^2).
\end{equation}
The complete expression is then
\begin{equation}
D=BX,\quad r=\operatorname{RowNorm}(D),\quad
U=\ones^{\mathsf T}u,\quad F=-B^{\mathsf T}\diag(w)D.
\label{eq:lj-force-tensor}
\end{equation}
\parallelmeaning Distances and analytic functions can be evaluated independently along the edge dimension, whereas endpoint contributions require reduction. This reorganizes the evaluation of the same LJ potential; it does not turn LJ into a neural network or a fixed matrix potential.

\subsubsection{Cutoffs require a fixed model definition}
The untruncated LJ potential has zero derivative at $r=2^{1/6}\sigma$. A force-shifted
truncation is defined by~\cite{lammps-shift}
\begin{equation}
\widetilde\phi(r)=
\begin{cases}
\phi(r)-\phi(r_c)-(r-r_c)\phi'(r_c),&r<r_c,\\
0,&r\ge r_c,
\end{cases}
\label{eq:force-shift}
\end{equation}
\begin{equation}
\widetilde\phi'(r)=
\begin{cases}\phi'(r)-\phi'(r_c),&r<r_c,\\0,&r\ge r_c.\end{cases}
\end{equation}
The edge operators can use these expressions in place of $\phi,\phi'$, but the conventional
reference implementation must use the same model. Force shifting changes forces inside the
cutoff and therefore changes the model, rather than merely rewriting its equations.

LJ is singular at $r_e=0$. Unused padding edges must be excluded safely before evaluating
$1/r_e$; multiplying an invalid result by zero does not reliably remove $0\times\infty$.
Parameter mixing, exclusions, and scaling weights are also part of the model specification.
For position-independent scaling coefficients $s_e$, the pair expression is
\begin{equation}
U_2=\sum_e s_e\phi_e(r_e),\qquad
g_e=s_e\frac{\phi'_e(r_e)}{r_e}D_e,\qquad F=-B^{\mathsf T}g.
\label{eq:pair-scaling}
\end{equation}
An excluded pair has $s_e=0$ and is skipped before singular functions are evaluated. The same
factor must multiply energy and its derivative. Type parameters are obtained from a specified
symmetric lookup or mixing rule, $\theta_e=\mathcal M(z_{i_e},z_{j_e};\Theta)$; no universal
mixing rule is assumed. For EAM, changes to density contributions and to pair energies must be
specified separately, since scaling a pair energy does not by itself scale an embedding term.
\subsection{Bonds, angles, and dihedrals on ordered hyperedges}
\label{sec:bonded}
\oldform For harmonic bonds and angles, respectively, we use
\begin{equation}
 U_b=\frac12\sum_b k_b(r_b-r_{0,b})^2,\qquad
 U_\theta=\frac12\sum_a k_a(\theta_a-\theta_{0,a})^2.
\end{equation}
In each angle $a=(i,j,k)$, the second atom is central. Let
$\vect a=\vect x_i-\vect x_j$ and $\vect b=\vect x_k-\vect x_j$, and define
$\vect g_a=\partial u_a/\partial\vect a$ and $\vect g_b=\partial u_a/\partial\vect b$. Then
\begin{equation}
 \vect f_i=-\vect g_a,\qquad
 \vect f_j=\vect g_a+\vect g_b,\qquad
 \vect f_k=-\vect g_b.
 \label{eq:angle-local-forces}
\end{equation}
\newform Harmonic bonds retain a two-endpoint incidence operator $B_b$:
\begin{equation}
 D_b=B_bX,\quad r_b=\operatorname{RowNorm}(D_b),\quad
 F_b=-B_b^{\mathsf T}\diag\!\left(\frac{k\odot(r_b-r_0)}{r_b}\right)D_b.
\end{equation}
For all angles, construct selection operators $S_i,S_j,S_k$ for the three atom roles. Geometry
and assembly become
\begin{align}
 A_\theta&=(S_i-S_j)X, & C_\theta&=(S_k-S_j)X,\\
 F_\theta&=-(S_i-S_j)^{\mathsf T}G_a-(S_k-S_j)^{\mathsf T}G_b.
 \label{eq:angle-tensor}
\end{align}
Here $G_a,G_b$ collect the local derivatives for each angle. Expanding
Eq.~\eqref{eq:angle-tensor} gives Eq.~\eqref{eq:angle-local-forces}; three-body geometry need
not be decomposed into pair potentials.

For a dihedral $\chi$ defined by four ordered atoms, with a potential such as
$u(\chi)=\kappa[1+\cos(n\chi-\delta)]$, the unified expression is
\begin{equation}
 U_\chi=\ones^{\mathsf T}u\!\left(\Gamma_\chi(\mathcal G_{I^{(4)}}X)\right),\qquad
 F_\chi=-\mathcal G_{I^{(4)}}^*\mathrm D\Gamma_\chi^*u'(\chi).
\end{equation}
The atom order, angular orientation, and image convention must be retained.
Appendix~\ref{app:bonded} gives the full conventional expressions, local gradients, and operator
counterparts for angles and dihedrals; Appendix~\ref{app:torsion} provides the reverse chain for
a dihedral. Comparisons with other software require parameter mapping. For example, the harmonic
bond and angle coefficients in LAMMPS already include the factor of
$1/2$~\cite{lammps-bond,lammps-angle}.
\subsection{Embedded-atom models with environment aggregation}
\label{sec:eam}
\subsubsection{Conventional single-element EAM}
The embedded-atom model (EAM) introduces many-body dependence through neighbor density and a
nonlinear embedding energy~\cite{daw1984,lammps-eam}. For a single element,
\begin{equation}
\rho_i=\sum_{j\in\mathcal N(i)}\psi(r_{ij}),\qquad
U=\sum_i\mathcal A(\rho_i)+\frac12\sum_i\sum_{j\in\mathcal N(i)}\phi(r_{ij}).
\label{eq:eam-traditional}
\end{equation}
An embedding-energy term cannot be treated as an independent pair potential. Differentiation
with respect to the distance of a pair $(i,j)$ gives
\begin{equation}
\frac{\partial U}{\partial r_{ij}}
=\phi'(r_{ij})+
\left[\mathcal A'(\rho_i)+\mathcal A'(\rho_j)\right]\psi'(r_{ij}).
\end{equation}
The conventional atomic force is therefore
\begin{equation}
\vect f_i=\sum_j
\left\{\phi'(r_{ij})+
[\mathcal A'(\rho_i)+\mathcal A'(\rho_j)]\psi'(r_{ij})\right\}
\frac{\vect x_j-\vect x_i}{r_{ij}}.
\label{eq:eam-traditional-force}
\end{equation}

\subsubsection{An operator expression with node--edge coupling}
Retain $B$ and $C=|B|$, with each unordered edge stored once. Since an edge contributes density
to both endpoints,
\begin{equation}
D=BX,\quad r=\operatorname{RowNorm}(D),\quad
\rho=C^{\mathsf T}\psi(r).
\label{eq:eam-density}
\end{equation}
The potential energy becomes
\begin{equation}
U=\ones_N^{\mathsf T}\mathcal A(\rho)+\ones_P^{\mathsf T}\phi(r).
\end{equation}
Let the node quantity be $b=\mathcal A'(\rho)$. The chain rule gives
\begin{align}
\dd U
&=b^{\mathsf T}\dd\rho+\phi'(r)^{\mathsf T}\dd r\\
&=b^{\mathsf T}C^{\mathsf T}\diag(\psi'(r))\dd r
+\phi'(r)^{\mathsf T}\dd r.\nonumber
\end{align}
Thus
\begin{equation}
w\equiv\nabla_rU=\phi'(r)+\psi'(r)\odot(Cb).
\label{eq:eam-w}
\end{equation}
The resulting force is
\begin{equation}
F=-B^{\mathsf T}\diag\!\left(\frac{w}{r}\right)D.
\label{eq:eam-force}
\end{equation}
Since $(Cb)_e=b_{i_e}+b_{j_e}$, expanding Eq.~\eqref{eq:eam-w} edge by edge gives
Eq.~\eqref{eq:eam-traditional-force}.


\paragraph{Dependencies in environment aggregation.}
$X\to D\to r\to\psi(r)\to\rho\to b\to w\to g\to F$.

Edges or nodes within each stage can be processed in parallel, but force evaluation must wait
for the density aggregation on which it depends. Tensor notation does not remove these many-body
dependencies.



\subsubsection{Multiple elements and asymmetric density contributions}
More generally, let $\psi_{\alpha\leftarrow\beta}(r)$ denote the density contributed by an atom
of type $\beta$ to a receiving atom of type $\alpha$. The conventional expression is
\begin{equation}
\rho_i=\sum_j\psi_{z_i\leftarrow z_j}(r_{ij}),\qquad
U=\sum_i\mathcal A_{z_i}(\rho_i)+\sum_{i<j}\phi_{z_i z_j}(r_{ij}).
\end{equation}
This form, in which both the receiver and contributor can affect the density, includes
EAM/FS-type cases~\cite{lammps-eam}. For edge $e=(i_e,j_e)$, define
\begin{equation}
p_e=\psi_{z_{i_e}\leftarrow z_{j_e}}(r_e),\quad
q_e=\psi_{z_{j_e}\leftarrow z_{i_e}}(r_e),\quad
b_i=\mathcal A'_{z_i}(\rho_i).
\end{equation}
The reformulation is
\begin{equation}
\rho=\src^{\mathsf T}p+\dst^{\mathsf T}q,\qquad
w=\phi'(r)+p'(r)\odot(\src b)+q'(r)\odot(\dst b).
\label{eq:eam-multi}
\end{equation}
Forces are again assembled using Eq.~\eqref{eq:eam-force}. If different terms have different
cutoffs, the candidate relation must cover the active range of every term used.
\section{Periodic geometry, neighbor validity, and time evolution}
\subsection{Discrete relations and periodic geometry}
\label{sec:pbc}
\subsubsection{Conventional and matrix expressions for periodic displacements}
For a general fixed cell, store the lattice vectors as rows of $\cell$. The conventional
minimum-image displacement can be defined by
\begin{equation}
\vect d_{ij}=\vect x_j-\vect x_i-\cell^{\mathsf T}\vect\lambda_{ij},\qquad
\vect\lambda_{ij}\in\argmin_{\vect n\in\ZZ^3}
\norm{\vect x_j-\vect x_i-\cell^{\mathsf T}\vect n}.
\end{equation}
\newform Collecting the integer image vectors into $\Lambda\in\ZZ^{P\times3}$ gives
\begin{equation}
D=BX-\Lambda\cell.
\label{eq:pbc-matrix}
\end{equation}
For local differentiation at fixed $\Lambda,\cell$, we have $\dd D=B\dd X$, so
\begin{equation}
F=-B^{\mathsf T}g
\end{equation}
still holds. With nonzero image offsets, however, it cannot be replaced by $F=-L_wX$ without an
additional term.

For a fixed orthorhombic box $\cell=\diag(L_x,L_y,L_z)$,
\begin{equation}
D_{e\alpha}=(BX)_{e\alpha}-L_\alpha
\round\!\left(\frac{(BX)_{e\alpha}}{L_\alpha}\right).
\end{equation}
For a general skew cell, component-wise rounding of fractional coordinates is not an
unconditionally valid nearest-lattice-point algorithm. The geometry operator must return a
minimizer of the stated distance problem and specify how ties are handled; the orthorhombic
rounding formula alone is not a general-cell implementation.

\paragraph{Minimum-image models and periodic image sums.}
Define the shortest nonzero lattice-translation length by
\begin{equation}
\ell_{\mathrm{lat}}=\min_{\vect n\in\ZZ^3\setminus\{0\}}
\norm{\cell^{\mathsf T}\vect n}.
\label{eq:lattice-length}
\end{equation}
For short-range pair or radial EAM terms that vanish for $r\ge r_c$, a sufficient condition
for at most one active image of each pair, and no self-image contributions, is
\begin{equation}
r_c<\frac12\ell_{\mathrm{lat}}.
\label{eq:unique-image-cutoff}
\end{equation}
Indeed, two images within $r_c$ of the same point would be separated by less than $2r_c$,
contradicting the definition of $\ell_{\mathrm{lat}}$. Under this condition the minimum-image
evaluation agrees with the corresponding finite-range periodic image sum. For an orthorhombic
box, $\ell_{\mathrm{lat}}=\min_\alpha L_\alpha$, giving the familiar half-box restriction
for short-range interactions~\cite{gromacs-pbc}. For a general cell, half the minimum
perpendicular face-to-face width is a possibly more restrictive sufficient bound; half the
shortest \emph{listed basis vector} is not valid for an arbitrary unreduced basis.

The physical cutoff $r_c$ and candidate radius $r_{\mathrm{list}}$ have different roles.
Requiring $r_{\mathrm{list}}<\ell_{\mathrm{lat}}/2$ is a convenient sufficient restriction
for a simple image-based list builder, but is not required by the local force identity.
A larger candidate radius is permissible for a builder that returns distinct unordered pairs,
covers all active interactions, and recomputes their nearest images. If the intended model
instead sums multiple active images, edges must include image labels and explicit counting
rules. Self-image terms require separate treatment; for a self-edge $B_e=0$, so $|B|$ cannot
represent its endpoint density contributions. Such models are outside the present simple-graph
formulation.

\paragraph{Regularity at image changes.}
For a smooth radial pair potential, the minimum distance, and hence its potential energy,
is continuous at an image tie away from collisions. Its gradient need not be continuous or
exist there. For example, along one periodic coordinate with period $L$,
\begin{equation}
U(s)=\phi\!\left(\min(s,L-s)\right),\quad 0<s<L,
\qquad
U'_-(L/2)=\phi'(L/2),\quad U'_+(L/2)=-\phi'(L/2).
\end{equation}
Thus an uncut LJ minimum-image energy generally has a cusp, rather than an energy jump.
Away from ties, its fixed-branch derivative is the actual local derivative. Condition
\eqref{eq:unique-image-cutoff} places image ties outside the active range and removes this
source of nonsmoothness. Regularity at the cutoff itself still depends on the truncation:
force shifting makes the radial energy continuously differentiable there, not generally twice
continuously differentiable. The periodic numerical example below tests only a fixed-branch
region of the uncut minimum-image model.

\subsubsection{Candidate and active interaction graphs}
\oldform A candidate list uses radius $r_{\mathrm{list}}=r_c+\delta$, where $\delta>0$ is the skin distance. The list need not be rebuilt at every step, but it must cover all currently active interactions~\cite{lammps-neighbor}.
\newform Treat graph construction as a discrete operator with an explicit contract:
\begin{equation}
\mathcal E_{\mathrm{cand}}=\mathcal R(X,\cell;r_c+\delta),\qquad
\mathcal E_{\mathrm{active}}(X)\subseteq\mathcal E_{\mathrm{cand}}.
\label{eq:neighbor-contract}
\end{equation}
The candidate graph may contain inactive edges, but it must not omit active ones. Cutoff
functions or validity masks determine the actual energy contributions.

\begin{proposition}[A sufficient reuse condition for a fixed box]
Let the maximum actual particle displacement since the last graph construction be
\begin{equation}
d_{\max}=\max_i\norm{\vect x_i(t)-\vect x_i(t_{\mathrm{build}})}.
\end{equation}
If displacements correctly account for periodic crossings and $2d_{\max}<\delta$, no pair that
moves from outside the candidate radius into the physical cutoff will be missed.
\end{proposition}
\begin{proof}
The change in relative displacement is bounded by the sum of the endpoint displacements, hence
by $2d_{\max}$. A pair initially farther apart than $r_c+\delta$ therefore remains farther apart
than $r_c$ under this condition. The periodic minimum distance satisfies the same bound,
provided that displacement tracking is consistent with particle motion rather than obtained by
directly subtracting wrapped coordinates.
\end{proof}
This condition does not directly cover cell deformation; NPT simulations must additionally
account for changes in the cell. Even when the candidate edges are reused, current distances and
images must remain correct. An outdated $\Lambda$ cannot be reused indefinitely.

\subsubsection{Differentiation and dynamic capacity}
Neighbor indices are discrete metadata. Position gradients are normally taken through continuous
geometry and energy on a currently valid candidate graph, rather than through sorting or graph
construction itself. This requires complete candidate coverage, consistent cutoff definitions,
and validity of the branch on which the derivative is evaluated.

If a fixed-capacity buffer overflows, the implementation must report the condition, increase
capacity, and recompute, rather than silently discard edges. JAX-MD likewise exposes
neighbor-buffer overflow explicitly~\cite{jax-neighbor}. Omitting interactions changes the
physical model and is a correctness failure, not an acceptable small numerical error.
\subsection{Velocity Verlet as a state map}
\label{sec:verlet}
\subsubsection{Conventional update equations}
Let $h=\Delta t$. The standard velocity Verlet scheme is~\cite{tuckerman1992,gromacs-md}
\begin{align}
\vect v_i^{n+1/2}&=\vect v_i^n+\frac{h}{2m_i}\vect f_i(X_n),\\
\vect x_i^{n+1}&=\vect x_i^n+h\vect v_i^{n+1/2},\\
\vect v_i^{n+1}&=\vect v_i^{n+1/2}+\frac{h}{2m_i}\vect f_i(X_{n+1}).
\label{eq:vv-traditional}
\end{align}
The second velocity update uses the potential-energy gradient at the new positions.

\subsubsection{Tensor update}
\begin{equation}
\begin{aligned}
V_{n+1/2}&=V_n+\frac h2\mass^{-1}F_n,\\
X_{n+1}&=X_n+hV_{n+1/2},\\
F_{n+1}&=-\nabla_XU(X_{n+1}),\\
V_{n+1}&=V_{n+1/2}+\frac h2\mass^{-1}F_{n+1}.
\end{aligned}
\label{eq:vv-tensor}
\end{equation}
After obtaining the new positions, candidate-graph validity must be established before
evaluating $F_{n+1}$. Initialization requires $F_0$; subsequent steps can reuse the force
obtained at the end of the preceding step.

\subsubsection{A symmetric composition of two elementary flows}
Define the kick and drift state maps by
\begin{align}
\mathcal K_h(X,V)&=(X,\,V+h\mass^{-1}F(X)),\\
\mathcal D_h(X,V)&=(X+hV,\,V).
\end{align}
Then
\begin{equation}
\mathcal T_h=\mathcal K_{h/2}\circ\mathcal D_h\circ\mathcal K_{h/2}.
\label{eq:vv-composition}
\end{equation}
Function composition is evaluated from right to left. Each application of $\mathcal K$ evaluates
the force at its own input positions, consistent with symmetric splitting of evolution
operators~\cite{tuckerman1992,gromacs-md}.

For sufficiently smooth conservative forces and under the usual stability conditions, the method
is second order in time. It does not conserve the original total energy exactly at every step.
Nonsmooth cutoffs, invalid neighbor lists, and inconsistent energy--force definitions can
compromise the expected numerical behavior.

\parallelmeaning Atom-wise kick and drift updates can be parallelized, but the causal dependencies between $X_{n+1}$, $F_{n+1}$, and $V_{n+1}$ must be retained. A hardware-independent operator reformulation does not make all time steps independently executable at once.
\section{Structural properties and implementation semantics}
\label{sec:validation}
\subsection{Representation invariance}

\textbf{Edge reversal.} Reversing a row $B_e$, together with its image vector in a periodic
system, changes the signs of both $D_e,g_e$ while leaving $u_e$ unchanged. The contribution
$-B_e^{\mathsf T}g_e$ is therefore unchanged. Any endpoint-dependent parameters or density
contributions must follow their endpoint roles.

\textbf{Edge permutation.} Permute all edge-associated parameters, images, and masks together.
For an edge permutation matrix $Q$, with $B'=QB$ and $g'=Qg$,
\begin{equation}
-(B')^{\mathsf T}g'=-B^{\mathsf T}Q^{\mathsf T}Qg=-B^{\mathsf T}g.
\end{equation}

\textbf{Atom relabeling.} For an atom permutation matrix $\Pi$, correctly permuting types and masses and transforming the indices gives $X'=\Pi X$ and $B'=B\Pi^{\mathsf T}$. Thus
\begin{equation}
U'=U,\qquad F'=\Pi F.
\end{equation}
Permuting coordinates without updating indices and types does not represent the same physical
system.

\subsection{Conditions for conservation laws}
For a potential depending on relative coordinates in the absence of external fields, translation
invariance gives
\begin{equation}
\sum_i\vect f_i=0.
\end{equation}
For an isolated nonperiodic system whose energy is also invariant under global rotations,
\begin{equation}
\sum_i\vect x_i\times\vect f_i=0.
\end{equation}
This torque identity should not be tested on an arbitrary periodic system by inserting wrapped
coordinates. Rotational tests must rotate the relevant cell and external fields as well, or be
restricted to nonperiodic systems without external fields.

Continuous conservative dynamics satisfies
\begin{equation}
\frac{\dd\mathcal H}{\dd t}
=\ip{\mass V}{\dot V}_F+\ip{\nabla_XU}{\dot X}_F
=\ip{V}{F}_F+\ip{-F}{V}_F=0.
\end{equation}
This is a property of the continuous equations; it does not imply exact energy conservation at
every discrete Verlet step.


\subsection{Operator interfaces and semantics}
\label{sec:operators}
\begin{table}[htbp]
\centering
\caption{Forward definitions and adjoint or validity semantics of the core operators. Overbars denote reverse-mode variables.}
\label{tab:operators}
\small
\begin{tabularx}{\textwidth}{@{}p{36mm}YY@{}}
\toprule
Operator & Forward definition & Reverse or correctness semantics\\
\midrule
\texttt{Gather} & $Q=\mathcal G_I(X)$ & $\overline X=\mathcal G_I^*(\overline Q)$\\
\texttt{Difference} & $D=BX-\Lambda\cell$ & $\overline X=B^{\mathsf T}\overline D$ at fixed indices, images, and cell\\
\texttt{LocalMap} & $Y_a=f_a(Q_a)$ & $\overline Q_a=\mathrm D f_a(Q_a)^*\overline Y_a$\\
\texttt{SegmentSum} & $Y_i=\sum_{e:s_e=i}Z_e$ & $\overline Z_e=\overline Y_{s_e}$\\
\texttt{ReduceSum} & $u=\sum_a y_a$ & $\overline y_a=\overline u$\\
\texttt{AssembleSum} & $F=\mathcal G_I^*(Z)$ & The reverse action on $Z$ is Gather\\
\texttt{NeighborRelation} & Candidate indices, capacity, validity & Discrete metadata; must cover all active edges\\
\texttt{StateUpdate} & State maps such as kick/drift & Preserve input--output dependencies and step-size semantics\\
\bottomrule
\end{tabularx}
\end{table}
The overbars in Table~\ref{tab:operators} indicate adjoint variables in reverse propagation, not
time averages. When several downstream operators depend on the same input, their reverse
contributions must also be summed.

An operator specification should include at least input/output shapes, continuous and discrete
parameters, units, boundary and degenerate cases, differentiation rules, validity masks, and
error conventions. Explicit construction of the full Jacobian is unnecessary.

\subsection{Mathematical tensors and physical intermediate buffers}
The expression
\begin{equation}
D\to r\to u,g\to U,F
\end{equation}
does not require separate allocations for $D,r,u,g$. A valid implementation may compute
quantities on demand in registers or local storage, fuse adjacent operators, and write only the
required outputs.

The same action $B^{\mathsf T}g$ can be implemented by edge-parallel evaluation with atomic
accumulation, atom-grouped reduction, local collection from a bidirectional full list, or other
equivalent strategies. Operator semantics should not prescribe a particular thread, warp, atomic
instruction, or memory layout.

\begin{equation}
\text{Physical semantics}\ \longrightarrow\
\text{Reordering/fusion}\ \longrightarrow\
\text{Backend execution}.
\end{equation}
Hardware independence makes the upper-level mathematical definitions reusable; it does not
guarantee that one schedule is optimal on every architecture.

\subsection{Floating-point and distributed semantics}
Floating-point addition is not strictly associative. Reduction trees, fused operations, and
precision choices can change rounding results~\cite{nvidia-fp}. It is therefore necessary to
distinguish mathematical equivalence in real arithmetic, agreement within a specified tolerance,
and bitwise reproducibility.

When interaction inputs are distributed across devices, Gather may require remote node data, and
its adjoint assembly may require returning or reducing contributions to their owners. A
graph--tensor representation retains these cross-partition dependencies and makes them part of
the operator contract.
\section{Small-system numerical verification}
\label{sec:experiments}
\subsection{Setup and error measures}
Verification uses the Python/NumPy implementation in \texttt{reproducibility/}, alongside
the manuscript source. It is independent of the repository's simulation engines: it imports
none of their force or integration routines. Within the script, explicit incidence/Gather
operators are compared with separately written scalar-loop energies and forces. The loop
dihedral uses a closed-form Cartesian gradient and a scalar-triple-product angle definition,
while the operator path uses the reverse chain in Appendix~\ref{app:torsion}. Both paths
implement the same stated model conventions; this is not a third-party software comparison.

The recorded run uses Python 3.9.6, NumPy 2.0.2, and double precision on an arm64 CPU.
The nonperiodic examples contain five particles. A periodic example contains four particles
in a fixed orthorhombic box. Coordinates, parameters, units, lists, velocities, and random
inputs are specified in Appendix~\ref{app:verification} and \texttt{config.json}.
The seven main tests include uncut LJ with $\epsilon=0.8$, $\sigma=0.9$ on all unordered
pairs. The periodic LJ example uses the closest image of each pair and stays away from image
ties; it does not approximate an infinite uncut periodic lattice sum. The EAM tests also use
all unordered pairs, with constructed smooth functions that test the embedding chain rule
rather than fitted material potentials.

The script and its recorded outputs were newly prepared for this revision. They replace the
earlier draft's numerical entries, whose original script and run metadata were unavailable;
they are not presented as recovered original-run data.

For each energy function, central differences are computed as
\begin{equation}
 F^{\rm FD}_{i\alpha}=-\frac{U(X+\eta E^{i\alpha})-U(X-\eta E^{i\alpha})}{2\eta},\qquad
 e_F=\max_{i,\alpha}|F_{i\alpha}-F^{\rm FD}_{i\alpha}|,
 \label{eq:finite-difference}
\end{equation}
where $E^{i\alpha}$ is $1$ at $(i,\alpha)$ and zero elsewhere. The finite-difference energy
is evaluated by the scalar-loop path; $F$ is obtained from the operator path. We also report
the dimensionless, scale-normalized error
\begin{equation}
e_{\mathrm{norm}}=\frac{e_F}{\max(F_*,\max_{i,\alpha}|F_{i\alpha}|)},\qquad
F_*=E_*/L_*,
\label{eq:normalized-force-error}
\end{equation}
where $E_*,L_*$ are the reference energy and length units. This is not a component-wise
relative error, which is ill-conditioned for nearly zero force components.

The script evaluates $\eta/L_*=10^{-4},10^{-5},10^{-6},10^{-7}$ and records every result;
Table~\ref{tab:numerics} uses $\eta/L_*=10^{-6}$, without selecting a different optimal
step for each model. These increments are distinct from the integration time step.
Central differences balance truncation and floating-point cancellation, so their errors need
not decrease monotonically as $\eta$ decreases. The configurations and stencils avoid
coincident particles, collinearity, and image ties. The checks require
$e_{\mathrm{norm}}\le10^{-7}$ at the tabulated step. This tolerance is an executable
consistency criterion for the chosen fixtures, not a bound on arbitrary configurations.

\subsection{Results}
Table~\ref{tab:numerics} reports the maximum absolute and scale-normalized force differences
in the new reference run. All seven absolute errors are below $7\times10^{-10}$ in the
reference force unit. These results support local derivative consistency for the listed
configurations, not an exhaustive proof for arbitrary inputs.
\begin{table}[tb]
\centering
\caption{Energy--force checks from the new double-precision reference run, at $\eta/L_*=10^{-6}$. The second column is in units of $E_*/L_*$; the third uses Eq.~\eqref{eq:normalized-force-error}. ``Periodic dihedral'' denotes angular periodicity; its coordinates are nonperiodic. The periodic LJ row is a local minimum-image check.}
\label{tab:numerics}
\begin{tabular}{@{}lrr@{}}
\toprule
Test model & $e_F/(E_*/L_*)$ & $e_{\mathrm{norm}}$\\
\midrule
% BEGIN GENERATED VERIFICATION ROWS
Lennard--Jones & $4.7743\times10^{-10}$ & $2.5400\times10^{-10}$\\
Harmonic bonds & $6.5197\times10^{-10}$ & $4.0346\times10^{-10}$\\
Harmonic angles & $5.7914\times10^{-10}$ & $4.9471\times10^{-10}$\\
Periodic dihedral potential & $5.5923\times10^{-10}$ & $2.4012\times10^{-10}$\\
Single-element EAM form & $3.0318\times10^{-10}$ & $3.0318\times10^{-10}$\\
Multi-element EAM form & $3.7896\times10^{-10}$ & $3.7896\times10^{-10}$\\
LJ in a fixed orthorhombic box & $1.9740\times10^{-10}$ & $1.9740\times10^{-10}$\\
% END GENERATED VERIFICATION ROWS
\bottomrule
\end{tabular}
\end{table}

Across the seven models, the maximum scalar-loop versus operator force difference was
$8.88\times10^{-16}$ in the reference force unit, including explicit bonded and dihedral
comparisons. For LJ, the half-list versus full-list force difference was
$2.22\times10^{-16}$ for both adjoint and source-only full-list assembly. The
extraction--assembly adjoint residual was $8.88\times10^{-16}$ in the numerical inner products.
One velocity Verlet step gave maximum position and velocity differences of $0$ and
$3.47\times10^{-18}$ in their respective reference units. These small residuals, including
the exact zero, are consistency checks for this run; they do not demonstrate bitwise agreement
on other hardware or independent validation of the physical conventions.

The archived JSON additionally records edge and atom permutation checks, nonperiodic net force
and torque, LJ virial strain derivatives, and a force-shifted LJ example with fixed exclusions
and scaling. It includes an excluded coincident pair, values at and approaching the cutoff,
and illustrative image-tie and skew-cell checks. These supplements do not constitute a general
geometry solver or a simulation-level validation suite.

\subsection{Scope of verification}
The checks cover local gradients, accumulation signs, counting factors, selected static cutoff
and geometry cases, and one state update. They do not test dynamic neighbor construction,
buffer overflow, trajectories crossing cutoff or image boundaries, long-term energy behavior,
variable-cell dynamics, or communication across devices. Finite differences can reveal a
mismatch between an energy and its derivative, but cannot exclude a convention error shared
by both. No external MD package was used in this reference run. GPU throughput, memory usage,
and scalability were not measured, so no speedup is reported. The complete inputs, script,
outputs, and environment metadata accompany the source in \texttt{reproducibility/}.
\section{Discussion and limitations}
\label{sec:extensions}
\subsection{Long-range electrostatics requires field or global operators}
The conventional Coulomb expression is
\begin{equation}
U_C=\frac12\sum_{i\ne j}\frac{q_iq_j}{4\pi\varepsilon_0r_{ij}}.
\end{equation}
The periodic long-range problem cannot be represented by a short-range neighbor graph
alone. Ewald and particle-mesh Ewald (PME) methods include real-space and reciprocal-space
contributions, together with self-interaction, boundary, and other
corrections~\cite{gromacs-electro}.

For \emph{a field-energy term in a selected discrete mesh model}, one may define
\begin{equation}
\rho_g=\mathcal S_Xq,\qquad
\varphi_g=\mathcal L^{-1}\rho_g,\qquad
U_g=\frac12\rho_g^{\mathsf T}\varphi_g.
\end{equation}
If $\mathcal L^{-1}$ is a fixed discrete operator symmetric under the chosen inner product, then
\begin{equation}
F_g=-\mathrm D_X(\mathcal S_Xq)^*\varphi_g.
\end{equation}
An FFT may implement $\mathcal L^{-1}$. This expression is not a complete PME derivation: charge
neutrality, the zero mode, mesh weights, and other corrections must be specified separately.

\subsection{Constraints require coupled solves}
Position constraints $c(X)=0$ and velocity constraints $J_c(X)v=0$ introduce coupled relations.
For a flattened velocity $v\in\RR^{3N}$ and mass matrix $\mathbb M=\mass\otimes I_3$, the
mass-metric velocity projection at fixed positions can be written as
\begin{equation}
v=v^*-\mathbb M^{-1}J^{\mathsf T}\lambda,\qquad
(J\mathbb M^{-1}J^{\mathsf T})\lambda=Jv^*.
\end{equation}
This expression requires appropriate treatment of constraint rank and solvability; nonlinear
position constraints additionally require iterative solution. We therefore reserve a
\texttt{Solve} semantic operation rather than assuming that independent \texttt{Map} evaluations
can handle all constraints.

\subsection{Machine-learning potentials share interfaces and derivative chains}
Classical potentials can be composed of analytic local functions and aggregations, while
energy-based machine-learning potentials use learnable maps of the form
\begin{equation}
U_\Theta=\mathcal N_\Theta(X,z,I),\qquad F=-\nabla_XU_\Theta.
\end{equation}
The common elements are state, indices, energy gradients, and assembly contracts. This does not
require all model internals to use the same dense matrix multiplication, nor does it eliminate
dependencies in multilayer message passing.

\section{Conclusions}
We have organized conventional indexed expressions for classical MD into extraction, geometry,
local transformations, aggregation, and adjoint assembly. When the energy decomposes into local
interaction terms, its energy and force expressions are
\begin{align}
 U(X)&=\sum_c\Phi_c\!\left(\Gamma_c(\mathcal G_cX);\Theta_c\right),\label{eq:master-energy}\\
 F(X)&=-\sum_c\mathcal G_c^*\mathrm D\Gamma_c(\mathcal G_cX)^*\nabla\Phi_c.
 \label{eq:master-force}
\end{align}
Here $c$ indexes interaction classes, and each $\Phi_c$ includes the sum over its local terms.
For models such as EAM that contain intermediate aggregations, differentiation follows the
dependent computational chain rather than treating embedding energies as independent pair
functions. Periodic images, interaction counting, and neighbor validity form part of the
mathematical contract alongside the operators. Time integration is organized through state maps
that preserve causal dependencies.

Under the stated assumptions, the conventional and reformulated expressions have the same
physical meaning. Small-system checks support local numerical consistency, but do not establish
stability or performance in production simulations. Subsequent implementations can use these
mathematical semantics as a reference layer while exploring sparse storage, operator fusion,
reduction strategies, and hardware scheduling, with correctness and performance evaluated
separately and the model definition held fixed.

\paragraph{Code and data availability.}
The \texttt{reproducibility/} directory beside this \LaTeX{} source contains
\texttt{verify\_equivalence.py}, \texttt{config.json}, \texttt{requirements.txt},
\texttt{verification\_results.json}, \texttt{environment.json}, a generated table-row fragment,
and a README with execution instructions and provenance. The JSON contains all tested force
arrays, finite-difference sweeps, and scalar residuals; source and configuration hashes connect
it to the recorded run. These files document the new reference calculation for this revision.
No archived software release or arXiv identifier is specified in this draft.

\begin{thebibliography}{99}
\small

\bibitem{thompson2022}
A.~P. Thompson, H.~M. Aktulga, R.~Berger, et al.
LAMMPS---a flexible simulation tool for particle-based materials modeling at the atomic, meso,
and continuum scales.
\emph{Computer Physics Communications}, 271:108171, 2022.
\href{https://doi.org/10.1016/j.cpc.2021.108171}{doi:10.1016/j.cpc.2021.108171}.

\bibitem{anderson2020}
J.~A. Anderson, J.~Glaser, and S.~C. Glotzer.
HOOMD-blue: A Python package for high-performance molecular dynamics and hard particle Monte
Carlo simulations.
\emph{Computational Materials Science}, 173:109363, 2020.
\href{https://doi.org/10.1016/j.commatsci.2019.109363}{doi:10.1016/j.commatsci.2019.109363}.
\href{https://arxiv.org/abs/1308.5587}{arXiv:1308.5587}.

\bibitem{edwards2014}
H.~C. Edwards, C.~R. Trott, and D.~Sunderland.
Kokkos: Enabling manycore performance portability through polymorphic memory access patterns.
\emph{Journal of Parallel and Distributed Computing}, 74(12):3202--3216, 2014.
\href{https://doi.org/10.1016/j.jpdc.2014.07.003}{doi:10.1016/j.jpdc.2014.07.003}.

\bibitem{schoenholz2020}
S.~S. Schoenholz and E.~D. Cubuk.
JAX, M.D.: A framework for differentiable physics.
\emph{Advances in Neural Information Processing Systems}, 33, 2020.
\href{https://arxiv.org/abs/1912.04232}{arXiv:1912.04232}.

\bibitem{daw1984}
M.~S. Daw and M.~I. Baskes.
Embedded-atom method: Derivation and application to impurities, surfaces, and other defects in metals.
\emph{Physical Review B}, 29:6443--6453, 1984.
\href{https://doi.org/10.1103/PhysRevB.29.6443}{doi:10.1103/PhysRevB.29.6443}.

\bibitem{blondel1996}
A. Blondel and M. Karplus.
New formulation for derivatives of torsion angles and improper torsion angles in molecular
mechanics: Elimination of singularities.
\emph{Journal of Computational Chemistry}, 17(9):1132--1141, 1996.
\href{https://doi.org/10.1002/(SICI)1096-987X(19960715)17:9<1132::AID-JCC5>3.0.CO;2-T}{Publisher DOI}.

\bibitem{tuckerman1992}
M. Tuckerman, B.~J. Berne, and G.~J. Martyna.
Reversible multiple time scale molecular dynamics.
\emph{The Journal of Chemical Physics}, 97(3):1990--2001, 1992.
\href{https://doi.org/10.1063/1.463137}{doi:10.1063/1.463137}.

\bibitem{thompson2009}
A.~P. Thompson, S.~J. Plimpton, and W. Mattson.
General formulation of pressure and stress tensor for arbitrary many-body interaction
potentials under periodic boundary conditions.
\emph{The Journal of Chemical Physics}, 131:154107, 2009.
\href{https://doi.org/10.1063/1.3245303}{doi:10.1063/1.3245303}.

\bibitem{jax-smap}
JAX-MD documentation. \emph{Higher Order Functions}.
\url{https://jax-md.readthedocs.io/en/main/jax_md.smap.html}.
Accessed September 17, 2026.

\bibitem{gromacs-pbc}
GROMACS Reference Manual. \emph{Periodic boundary conditions}.
\url{https://manual.gromacs.org/current/reference-manual/algorithms/periodic-boundary-conditions.html}.
Accessed September 17, 2026.

\bibitem{lammps-lj}
LAMMPS documentation. \emph{pair\_style lj/cut command}.
\url{https://docs.lammps.org/pair_lj.html}.
Accessed September 16, 2026.

\bibitem{lammps-eam}
LAMMPS documentation. \emph{pair\_style eam command}.
\url{https://docs.lammps.org/pair_eam.html}.
Accessed September 16, 2026.

\bibitem{lammps-neighbor}
LAMMPS Programmer Guide. \emph{4.4.3 Neighbor lists}.
\url{https://docs.lammps.org/Developer_par_neigh.html}.
Accessed September 16, 2026.

\bibitem{lammps-bond}
LAMMPS documentation. \emph{bond\_style harmonic command}.
\url{https://docs.lammps.org/bond_harmonic.html}.
Accessed September 16, 2026.

\bibitem{lammps-angle}
LAMMPS documentation. \emph{angle\_style harmonic command}.
\url{https://docs.lammps.org/angle_harmonic.html}.
Accessed September 16, 2026.

\bibitem{lammps-shift}
LAMMPS documentation. \emph{pair\_style lj/smooth/linear command}.
\url{https://docs.lammps.org/pair_lj_smooth_linear.html}.
Accessed September 16, 2026.

\bibitem{gromacs-md}
GROMACS Reference Manual. \emph{Molecular Dynamics}.
\url{https://manual.gromacs.org/current/reference-manual/algorithms/molecular-dynamics.html}.
Accessed September 16, 2026.

\bibitem{jax-neighbor}
JAX-MD documentation. \emph{Spatial Partitioning}.
\url{https://jax-md.readthedocs.io/en/main/jax_md.partition.html}.
Accessed September 16, 2026.

\bibitem{nvidia-fp}
NVIDIA. \emph{Floating Point and IEEE 754}.
\url{https://docs.nvidia.com/cuda/floating-point/index.html}.
Accessed September 16, 2026.

\bibitem{gromacs-electro}
GROMACS Reference Manual. \emph{Long Range Electrostatics}.
\url{https://manual.gromacs.org/current/reference-manual/functions/long-range-electrostatics.html}.
Accessed September 16, 2026.
\end{thebibliography}
\appendix
\numberwithin{equation}{section}
\section{Detailed formula comparisons for bonded interactions}
\label{app:bonded}
\subsection{Harmonic bonds}
\oldform For a bond set $\mathcal B$, define
\begin{equation}
U_b=\frac12\sum_{b\in\mathcal B}k_b(r_b-r_{0,b})^2,
\qquad
\frac{\dd u_b}{\dd r_b}=k_b(r_b-r_{0,b}).
\end{equation}
\newform Replace the neighbor-edge incidence operator with a bond incidence operator $B_b$:
\begin{equation}
D_b=B_bX,\qquad r_b=\operatorname{RowNorm}(D_b),
\end{equation}
\begin{equation}
U_b=\frac12\ones^{\mathsf T}\!\left[k\odot(r_b-r_0)^2\right],\qquad
F_b=-B_b^{\mathsf T}\diag\!\left(\frac{k\odot(r_b-r_0)}{r_b}\right)D_b.
\end{equation}
The LAMMPS harmonic bond energy is written as $K(r-r_0)^2$, with the factor of $1/2$ absorbed
into the coefficient~\cite{lammps-bond}. Comparison with the present convention therefore
requires $k=2K$, rather than copying parameter values directly.

Bond topology is generally independent of the nonbonded cutoff. Under periodic boundaries, bond
images must also be consistent with molecular topology; arbitrary, mutually inconsistent
coordinate copies cannot be assigned to a complex molecule crossing the box.

\subsection{Harmonic angles: the full three-atom gradient}
\oldform For an angle $a=(i,j,k)$ centered at $j$, define
\begin{equation}
\vect a=\vect x_i-\vect x_j,\qquad
\vect b=\vect x_k-\vect x_j,\qquad
\alpha=\norm{\vect a},\quad\beta=\norm{\vect b},
\end{equation}
\begin{equation}
c=\frac{\vect a\cdot\vect b}{\alpha\beta},\qquad
\theta=\arccos c,\qquad
u_a=\frac12k_a(\theta-\theta_{0,a})^2.
\end{equation}
In the nondegenerate region $\alpha,\beta>0$ and $|c|<1$,
\begin{align}
\frac{\partial\theta}{\partial\vect a}
&=-\frac1{\sqrt{1-c^2}}
\left(\frac{\vect b}{\alpha\beta}-\frac{c\vect a}{\alpha^2}\right),\\
\frac{\partial\theta}{\partial\vect b}
&=-\frac1{\sqrt{1-c^2}}
\left(\frac{\vect a}{\alpha\beta}-\frac{c\vect b}{\beta^2}\right).
\end{align}
Let
\begin{equation}
\vect g_a=k_a(\theta-\theta_{0,a})\frac{\partial\theta}{\partial\vect a},
\qquad
\vect g_b=k_a(\theta-\theta_{0,a})\frac{\partial\theta}{\partial\vect b}.
\end{equation}
The local atomic forces are then
\begin{equation}
\vect f_i=-\vect g_a,\qquad
\vect f_j=\vect g_a+\vect g_b,\qquad
\vect f_k=-\vect g_b.
\label{eq:angle-local-forces-app}
\end{equation}

\newform Construct three selection operators $S_i,S_j,S_k$ for the angle set. Let
\begin{equation}
A_\theta=(S_i-S_j)X,\qquad
C_\theta=(S_k-S_j)X.
\end{equation}
Evaluate the angles and derivatives above row by row, collecting them into
$G_a,G_b\in\RR^{P_3\times3}$. Then
\begin{equation}
F_\theta=-(S_i-S_j)^{\mathsf T}G_a
-(S_k-S_j)^{\mathsf T}G_b.
\label{eq:angle-tensor-app}
\end{equation}
Equivalently, using the unified hyperedge extraction notation,
\begin{equation}
U_\theta=\Phi_\theta\circ\Gamma_\theta\circ\mathcal G_{I^{(3)}}(X),
\qquad
F_\theta=-\mathcal G_{I^{(3)}}^*\,
\mathrm D\Gamma_\theta^*\,\nabla\Phi_\theta.
\end{equation}
LAMMPS harmonic angles also use the coefficient convention
$K(\theta-\theta_0)^2$~\cite{lammps-angle}. Angles in the present calculations are in radians;
parameter-unit conversion and collinear degeneracies require explicit treatment.

\subsection{Periodic dihedrals and orientation conventions}
\oldform For an ordered quadruplet $(i,j,k,l)$, choose the definitions
\begin{gather}
\vect b_1=\vect x_j-\vect x_i,\quad
\vect b_2=\vect x_k-\vect x_j,\quad
\vect b_3=\vect x_l-\vect x_k,\\
\vect n_1=\vect b_1\times\vect b_2,\quad
\vect n_2=\vect b_2\times\vect b_3,\quad
\widehat{\vect b}_2=\vect b_2/\norm{\vect b_2},\\
x=\vect n_1\cdot\vect n_2,\qquad
y=\widehat{\vect b}_2\cdot(\vect n_1\times\vect n_2),\qquad
\chi=\atan(y,x).
\label{eq:dihedral-definition}
\end{gather}
For illustration, define the periodic potential
\begin{equation}
u(\chi)=\kappa[1+\cos(n\chi-\delta)],\qquad
u'(\chi)=-\kappa n\sin(n\chi-\delta),
\end{equation}
where $n$ is an integer. The conventional forces are
\begin{equation}
\vect f_p=-u'(\chi)\frac{\partial\chi}{\partial\vect x_p},
\qquad p\in\{i,j,k,l\}.
\end{equation}
\newform With $Q=\mathcal G_{I^{(4)}}(X)$ and $\chi=\Gamma_\chi(Q)$,
\begin{equation}
U_\chi=\ones^{\mathsf T}u(\chi),\qquad
F_\chi=-\mathcal G_{I^{(4)}}^*\,
\mathrm D\Gamma_\chi(Q)^*\,u'(\chi).
\label{eq:dihedral-reform}
\end{equation}
Torsion derivatives have established molecular-mechanics formulations~\cite{blondel1996}.
For the convention specified here, local differentiation requires only successive applications of dot products, cross products,
normalization, and
\begin{equation}
\dd\chi=\frac{x\,\dd y-y\,\dd x}{x^2+y^2}.
\label{eq:atan-derivative}
\end{equation}
Appendix~\ref{app:torsion} gives the reverse chain without constructing a full Jacobian.

Software packages may use different atom orders and phase conventions, which must be mapped
before comparison. Cases with $\norm{\vect b_2}=0$ or degenerate normal vectors require explicit
treatment. A jump in the dihedral angle branch does not by itself imply a discontinuity in the
periodic energy.
\section{Reverse chain for the local dihedral gradient}
\label{app:torsion}
This appendix uses the \emph{specific orientation convention} of
Eq.~\eqref{eq:dihedral-definition}. For brevity, let $\vect t=\widehat{\vect b}_2$, $\ell=\norm{\vect
b_2}$, and $c_\chi=u'(\chi)$. First,
\begin{equation}
\overline x=-\frac{c_\chi y}{x^2+y^2},\qquad
\overline y=\frac{c_\chi x}{x^2+y^2}.
\end{equation}
The reverse rules for the dot and cross products give
\begin{align}
\overline{\vect n}_1&=\overline x\,\vect n_2
+\overline y\,(\vect n_2\times\vect t),\\
\overline{\vect n}_2&=\overline x\,\vect n_1
+\overline y\,(\vect t\times\vect n_1),\\
\overline{\vect t}&=\overline y\,(\vect n_1\times\vect n_2).
\end{align}
Normalization and the two cross products yield
\begin{align}
\overline{\vect b}_1&=\vect b_2\times\overline{\vect n}_1,\\
\overline{\vect b}_2&=
\frac{\overline{\vect t}-\vect t(\vect t\cdot\overline{\vect t})}{\ell}
+\overline{\vect n}_1\times\vect b_1
+\vect b_3\times\overline{\vect n}_2,\\
\overline{\vect b}_3&=\overline{\vect n}_2\times\vect b_2.
\end{align}
Finally, the local coordinate gradients are
\begin{equation}
\begin{aligned}
\nabla_{\vect x_i}u&=-\overline{\vect b}_1,\qquad &
\nabla_{\vect x_j}u&=\overline{\vect b}_1-\overline{\vect b}_2,\\
\nabla_{\vect x_k}u&=\overline{\vect b}_2-\overline{\vect b}_3, &
\nabla_{\vect x_l}u&=\overline{\vect b}_3.
\end{aligned}
\end{equation}
Negating these gradients and accumulating them through $\mathcal G_{I^{(4)}}^*$ gives the force
in Eq.~\eqref{eq:dihedral-reform}. The calculation uses only local vector operations and does
not require a large dense Jacobian.
\section{Tensor expressions for energy, temperature, and configurational virial}
\label{sec:observables}
\subsection{Kinetic energy and temperature}
\oldform
\begin{equation}
K=\frac12\sum_i m_i\vect v_i\cdot\vect v_i,\qquad
T=\frac{2K}{N_{\mathrm{dof}}k_B}.
\end{equation}
\newform
\begin{equation}
K=\frac12\tr(V^{\mathsf T}\mass V),\qquad
T=\frac{\tr(V^{\mathsf T}\mass V)}{N_{\mathrm{dof}}k_B},\qquad
\mathcal H=K+U.
\end{equation}
The degrees of freedom must be consistent with constraints and center-of-mass removal.
Temperature must use velocities at the time point specified by the integration
convention~\cite{gromacs-md}.

\subsection{Virial for pair potentials and radial EAM}
To avoid ambiguity in signs and factors of $1/2$ across software packages, we \emph{define} the
configurational virial as
\begin{equation}
W_{\alpha\beta}=-\sum_{e=1}^P D_{e\alpha}g_{e\beta}.
\label{eq:virial-component}
\end{equation}
Under our orientation convention, $g_e$ is the edge force acting on the source endpoint $i_e$,
while $D_e$ points toward the target endpoint. The reformulated expression is
\begin{equation}
W=-D^{\mathsf T}g.
\label{eq:virial-tensor}
\end{equation}
It can also be defined from the energy derivative under the uniform deformation $D\mapsto
D(I+\varepsilon)$:
\begin{equation}
\delta U=\ip{D^{\mathsf T}g}{\delta\varepsilon}_F,
\qquad W=-\left.\nabla_\varepsilon U\right|_{\varepsilon=0}.
\end{equation}
For a periodic system containing only these short-range internal forces, with no constraints,
external fields, or other corrections, the pressure tensor under this convention is
\begin{equation}
\mathsf P=\frac{V^{\mathsf T}\mass V+W}{\Omega},\qquad
p=\frac13\tr(\mathsf P).
\end{equation}
GROMACS uses a different configurational virial symbol $\Xi$ and prefactor. Comparison requires
the mapping $W=-2\Xi$, rather than matching variable names alone~\cite{gromacs-md}.

For periodic systems, $X^{\mathsf T}F$ formed from wrapped coordinates cannot replace
Eq.~\eqref{eq:virial-tensor}, which uses the correct image displacements. Additional
contributions from angles, constraints, long-range terms, and explicit cell dependence must be
derived from the complete energy or constraint definitions.
For a broader treatment of periodic virials with many-body potentials, see
Thompson, Plimpton, and Mattson~\cite{thompson2009}.
\section{Algorithms at the semantic level}
\label{app:algorithms}
Algorithm~\ref{alg:pair} specifies the mathematical evaluation order for a general radial pair
potential, and Algorithm~\ref{alg:verlet} specifies one state update. A backend may fuse or
eliminate the named intermediates. The pseudocode prescribes neither a device-memory layout nor
a separate kernel for each line. The potential functions include the selected cutoff, exclusion,
and scaling rules.

\begin{algorithm}[htbp]
\caption{Pair energy and force evaluation with incidence operators}\label{alg:pair}
\begin{algorithmic}[1]
\Require $X$, valid candidate half list $I$, current images $\Lambda$, cell $\cell$, potential parameters $\Theta$
\Ensure $U,F$, and optional configurational virial $W$
\State Define $B=\dst-\src$ from $I$; only endpoint indices need to be stored
\State $D\gets BX-\Lambda\cell$; evaluate $r_e\gets\|D_e\|$ on valid edges
\State Evaluate $u_e\gets\phi_e(r_e;\Theta)$ and $g_e\gets\phi'_e(r_e;\Theta)D_e/r_e$
\State $U\gets\sum_e u_e$
\State $F\gets-B^{\mathsf T}g$ \Comment{Sum contributions at repeated atom indices}
\State Optionally, $W\gets-D^{\mathsf T}g$
\State \Return $U,F,W$
\end{algorithmic}
\end{algorithm}

\begin{algorithm}[htbp]
\caption{One velocity Verlet step with neighbor-validity checks}\label{alg:verlet}
\begin{algorithmic}[1]
\Require $X_n,V_n,F_n$, mass matrix $\mass$, step size $h$, cached candidate relation
\Ensure $X_{n+1},V_{n+1},F_{n+1}$
\State $V_{n+1/2}\gets V_n+(h/2)\mass^{-1}F_n$
\State $X_{n+1}\gets X_n+hV_{n+1/2}$
\If{the candidate relation no longer satisfies the validity condition}
\State Rebuild the relation at $X_{n+1}$; report overflow, expand capacity, and rebuild if needed
\EndIf
\State Update current images and continuous geometry; do not reuse stale displacements
\State $(U_{n+1},F_{n+1})\gets\operatorname{Potential}(X_{n+1},I,\Theta)$
\State $V_{n+1}\gets V_{n+1/2}+(h/2)\mass^{-1}F_{n+1}$
\State \Return $X_{n+1},V_{n+1},F_{n+1}$
\end{algorithmic}
\end{algorithm}

For many-body potentials, \texttt{Potential} must retain all intermediate aggregation
dependencies. Algorithm~\ref{alg:pair} cannot directly replace the EAM density--embedding--force
chain. In a multi-device implementation, this call must also obtain remote inputs and reduce
contributions to their owners.
\section{Reproducibility details for the reference implementation}
\label{app:verification}

The accompanying \texttt{reproducibility/verify\_equivalence.py} uses Python and NumPy with
\texttt{float64} arithmetic. All numerical values below are expressed in fixed reference units
$L_*$ for length, $E_*$ for energy, $M_*$ for mass, and a chosen density unit $\rho_*$ for
the constructed EAM functions. Time and force units are
$t_*=L_*\sqrt{M_*/E_*}$ and $F_*=E_*/L_*$. Angles are in radians. Exponential functions
below take the numerical distance $r/L_*$ as their argument; coefficients carry the units
needed by their respective energy and density definitions. No material-specific unit mapping
is asserted.

The nonperiodic tests fix the following five-particle configuration, displayed as $X/L_*$:
\begin{equation}
X=\begin{bmatrix}
0.10&0.20&0.00\\
1.40&0.10&0.20\\
0.60&1.45&-0.10\\
1.80&1.40&0.65\\
-0.80&0.70&0.50
\end{bmatrix}.
\end{equation}
LJ uses $\epsilon/E_*=0.8,\sigma/L_*=0.9$. LJ and both EAM examples use all ten unordered
pairs, in lexicographic order $(0,1),(0,2),\ldots,(3,4)$, with no cutoff or exclusions.
The single-element EAM test uses constructed smooth functions
solely to check the chain rule:
\begin{equation}
\psi(r)=e^{-r},\qquad \mathcal A(\rho)=0.4\rho^2,\qquad
\phi(r)=0.15e^{-2r}.
\end{equation}
These functions are \textbf{not fitted material potentials}, and the tests do not validate
material properties. The multi-element example uses asymmetric density amplitudes to check
that force differentiation preserves the receiver direction specified by its energy in
Eq.~\eqref{eq:eam-multi}; this is not an external validation of an EAM file convention.


\subsection{Additional parameters and execution}
The zero-based harmonic bond indices are $(0,1),(1,2),(2,3),(3,4)$, with
$k=(1.2,1.7,2.1,0.8)$, $r_0=(1.0,1.1,1.2,1.1)$.
The angle indices are $(0,1,2),(1,2,3),(2,3,4)$, with
$k=(1.1,1.3,0.9)$ and $\theta_0=(1.6,1.7,1.8)$ radians.
The dihedral indices are $(0,1,2,3),(1,2,3,4)$, with each term using
$\kappa=0.7$, $n=3$, $\delta=0.2$.
Mathematical atom labels in the text start at $1$, whereas script indices start at $0$.

The multi-element example has labels $z=(0,1,0,1,1)$ and density-amplitude matrix
\begin{equation}
 A_{\alpha\leftarrow\beta}=\begin{bmatrix}1.0&0.7\\1.3&0.9\end{bmatrix},\qquad
 \psi_{\alpha\leftarrow\beta}(r)=A_{\alpha\leftarrow\beta}e^{-r}.
\end{equation}
The embedding functions are $\mathcal A_\alpha(\rho)=a_\alpha\rho^2$ with $a=(0.3,0.5)$; the
pair term remains $0.15e^{-2r}$. These asymmetric parameters test the receiver/contributor
direction and are not material parameters.

The single integration step uses the nonperiodic uncut LJ model above,
$m/M_*=(1.0,1.3,0.9,1.2,1.1)$, $h/t_*=10^{-3}$, and
\begin{equation}
\frac{V_0}{L_*/t_*}=\begin{bmatrix}
0.10&-0.05&0.02\\
-0.03&0.04&0.01\\
0.02&0.01&-0.04\\
-0.01&-0.02&0.03\\
0.00&0.03&-0.02
\end{bmatrix}.
\end{equation}
The initial forces are evaluated separately by the scalar-loop and operator paths.

The periodic example uses all six unordered pairs of
\begin{equation}
\frac{X_{\mathrm{pbc}}}{L_*}=\begin{bmatrix}
0.2&0.3&0.1\\
3.2&0.5&3.0\\
1.4&3.0&1.3\\
2.6&1.7&0.3
\end{bmatrix},\qquad
\frac{\cell}{L_*}=\diag(4.0,4.2,3.8).
\end{equation}
It has no cutoff. The smallest component-wise distance of an image displacement to a
half-box boundary is $0.2L_*$, so all stated finite-difference stencils remain on the same
image branch. Images are nevertheless recomputed for each perturbed energy evaluation.

The adjoint test uses the zero-based index array
\begin{equation}
I=\begin{bmatrix}0&1&0\\2&1&3\\4&2&1\end{bmatrix}
\end{equation}
and a tensor $Z\in\RR^{3\times3\times3}$ generated by exactly one call to
\texttt{numpy.random.default\_rng(17).standard\_normal((3,3,3))}, with no preceding draws.
The repeated index in the first term tests additive accumulation of multiple slots. The
actual tensor is also saved in the JSON, so its values need not be inferred from a seed alone.

The supplementary LJ test applies force shifting with $r_c/L_*=1.6$ and scales
$s=(0,0.5,1,1,0.5,1,1,1,1,1)$ to the ten lexicographically ordered nonperiodic pairs.
Additional static checks skip an excluded coincident pair before evaluating LJ, evaluate
the cutoff value and one-sided approach, and differentiate the LJ virial with strain increment
$10^{-6}$. All fixtures and supplementary geometry examples are recorded in the JSON.

From the manuscript directory, the execution command is
\begin{quote}\ttfamily python reproducibility/verify\_equivalence.py\end{quote}
The script also runs from other working directories because default input and output paths
are resolved relative to its own file. It writes \texttt{verification\_results.json},
\texttt{environment.json}, and \texttt{table\_rows.tex} into \texttt{reproducibility/}; it
does not edit this manuscript. The optional \texttt{--output-dir} argument preserves the
archived run while creating another, and \texttt{--check-manuscript} compares the marked
table rows with the current run. Exact row matching is an editorial synchronization check,
not a portable numerical acceptance criterion.

The JSON records every finite-difference force array and residual; the environment file records
Python, NumPy, platform and numerical-library information, the execution time, and SHA-256
hashes of the script and configuration. The reference implementation deliberately constructs
small incidence matrices; it is not a sparse backend for large systems. Reproduction is
judged against the stated tolerances, and small last-digit changes across environments are
expected. The README gives dependency installation and execution details.

\end{document}
